%%
%% Copyright 2007-2025 Elsevier Ltd
%%
%% This file is part of the 'Elsarticle Bundle'.
%% ---------------------------------------------
%%
%% It may be distributed under the conditions of the LaTeX Project Public
%% License, either version 1.3 of this license or (at your option) any
%% later version.  The latest version of this license is in
%%    http://www.latex-project.org/lppl.txt
%% and version 1.3 or later is part of all distributions of LaTeX
%% version 1999/12/01 or later.
%%
%% The list of all files belonging to the 'Elsarticle Bundle' is
%% given in the file `manifest.txt'.

%% Template article for Elsevier's document class `elsarticle'
%% with harvard style bibliographic references
\documentclass[preprint,12pt,authoryear]{elsarticle}

%% Use the option review to obtain double line spacing
% \documentclass[authoryear,preprint,review,12pt]{elsarticle}

%% Use the options 1p,twocolumn; 3p; 3p,twocolumn; 5p; or 5p,twocolumn
%% for a journal layout:
% \documentclass[final,1p,times,authoryear]{elsarticle}
% \documentclass[final,1p,times,twocolumn,authoryear]{elsarticle}
% \documentclass[final,3p,times,authoryear]{elsarticle}
% \documentclass[final,3p,times,twocolumn,authoryear]{elsarticle}
% \documentclass[final,5p,times,authoryear]{elsarticle}
% \documentclass[final,5p,times,twocolumn,authoryear]{elsarticle}

\usepackage[english]{babel}
\usepackage{amssymb}
\usepackage{amsmath}
\usepackage{amsthm}

%% The lineno packages adds line numbers. Start line numbering with
%% \begin{linenumbers}, end it with \end{linenumbers}. Or switch it on
%% for the whole article with \linenumbers.
\usepackage{lineno}

\journal{Computer \& Graphics}

\begin{document}

\begin{frontmatter}

	%% Title, authors and addresses

	%% use the tnoteref command within \title for footnotes;
	%% use the tnotetext command for theassociated footnote;
	%% use the fnref command within \author or \affiliation for footnotes;
	%% use the fntext command for theassociated footnote;
	%% use the corref command within \author for corresponding author footnotes;
	%% use the cortext command for theassociated footnote;
	%% use the ead command for the email address,
	%% and the form \ead[url] for the home page:
	%% \title{Title\tnoteref{label1}}
	%% \tnotetext[label1]{}
	%% \author{Name\corref{cor1}\fnref{label2}}
	%% \ead{email address}
	%% \ead[url]{home page}
	%% \fntext[label2]{}
	%% \cortext[cor1]{}
	%% \affiliation{organization={},
	%%            addressline={},
	%%            city={},
	%%            postcode={},
	%%            state={},
	%%            country={}}
	%% \fntext[label3]{}

	\title{Rebuilding 3D models using rule-based geometric modeling operations}

	\author[lias]{Maxime Gaide}
	\author[lias]{David Marcheix}
	\author[xlim]{Agn\`{e}s Arnould}
	\author[xlim]{Xavier Skapin}
	\author[xlim]{Hakim Belhaouari}
	\author[lias]{St\'{e}phane Jean}
	% \cortext[cor1]{contacts}

	%% Author affiliation
	\affiliation[lias]{%
		cityorganization={University of Poitiers, ISAE-ENSMA, LIAS},
		% addressline={},
		city={Poitiers},
		postcode={86000},
		% state={},
		country={France}}%

	\affiliation[xlim]{%
		organization={University of Poitiers, Univ. Limoges, CNRS, XLIM},
		% addressline={},
		city={Poitiers},
		postcode={86000},
		% state={},
		country={France}}%

	%% Abstract
	\begin{abstract}
	  Regenerating a 3D parametric model is a tedious process of trial and error. %
	  In order to make the model design process easier, current parametric
		modeling systems propose a rebuild feature allowing users to edit the
		construction process of a model and regenerate it. %
	  However, most modeling systems require the tracking of numerous topological
		entities and only consider the editing from the standpoint of geometric
		parameters, thus limiting the rebuilding capabilities. %

	  In this paper, we present a rebuild mechanism allowing users to edit
		construction processes through the addition and deletion of operations while
		it tracks identified topological parameters only. %

		Our approach relies on the formalization of objects with generalized maps
		and modeling operations with graph transformation rules. %
		This make it possible to syntactically analyze operations and tracks
		topological changes (creation, deletion, merging,\dots) during the build and
		rebuild processes. %

	\end{abstract}

	%%Graphical abstract
	\begin{graphicalabstract}
		%\includegraphics{grabs}
	\end{graphicalabstract}

	%%Research highlights
	\begin{highlights}
		\item Research highlight 1
		\item Research highlight 2
	\end{highlights}

	%% Keywords
	\begin{keyword}
		%% keywords here, in the form: keyword \sep keyword

		%% PACS codes here, in the form: \PACS code \sep code

		%% MSC codes here, in the form: \MSC code \sep code
		%% or \MSC[2008] code \sep code (2000 is the default)

	\end{keyword}

\end{frontmatter}

\linenumbers

\section{Introduction}
\label{sec:introduction}

\section{Related works}
\label{sec:related-works}



\section{Event detection}
\label{sec:event-detection}

\subsection{Creation}
\label{sec:creation}

\subsection{Splitting}
\label{sec:splitting}


\subsection{Deletion}
\label{sec:deletion}


\subsection{Merging}
\label{sec:merging}

\subsection{Non-modification}
\label{sec:non-modification}

\subsection{Modification}
\label{sec:modification}



\section{Rule-based rebuilding}
\label{sec:rule-based-rebu}

\subsection{Persistent names for darts}
\label{sec:pers-names-darts}

\subsection{Persistent names for orbits}
\label{sec:pers-names-orbits}

\subsubsection{Reconstitution of orbits histories}
\label{sec:reconst-orbits-hist}

\subsubsection{Integration of origins}
\label{sec:integration-origins}

\subsection{Rebuild mechanisms}
\label{sec:rebuild-mechanisms}

\subsubsection{Orbit matching}
\label{sec:orbit-matching}





\bibliographystyle{elsarticle-harv}
\bibliography{references}
\end{document}

\endinput
%%
%% End of file `elsarticle-template-harv.tex'.
